% Corporate purchase order - Pressa public template (MIT). https://pressa.dev/templates/purchase-order?design=corporate-purchase-order
% Line totals, subtotal, tax and total are computed from the data, each rounded to the
% currency's minor unit (cents, whole yen) as printed, so the printed numbers add up.
% Unit prices keep up to 4 decimals for items priced below a cent.
\documentclass[10pt,a4paper]{article}
\usepackage[T1]{fontenc}
\usepackage[default,semibold]{sourcesanspro}
\usepackage[margin=18mm,top=24mm]{geometry}
\usepackage[table]{xcolor}
\usepackage{tabularx}
\usepackage{longtable}
\usepackage{array}
\usepackage{eso-pic}
\usepackage{tikz}
\usepackage[hidelinks]{hyperref}

% `accent` is the colour a visitor can change; tints mix it inline (accent!8),
% never through \colorlet, so they follow the change.
\definecolor{accent}{HTML}{374151}
\definecolor{ink}{HTML}{1F2937}
\definecolor{muted}{HTML}{6B7280}
\pagestyle{empty}
\setlength{\parindent}{0pt}
\newcolumntype{L}{>{\raggedright\arraybackslash}X}
\newcolumntype{R}{>{\raggedleft\arraybackslash}X}
\newcommand{\fieldname}[1]{\textcolor{muted}{\footnotesize\MakeUppercase{#1}}}
% A band of the accent colour across the top of every page.
\AddToShipoutPictureBG{\AtPageUpperLeft{\raisebox{-8mm}{\textcolor{accent}{\rule{\paperwidth}{8mm}}}}}

\begin{document}
\color{ink}
{% assign d = currency | currency_decimals %}{% assign subtotal = 0 %}{% for item in items %}{% assign price = item.unit_price | round: 4 %}{% assign line = item.quantity | times: price | round: d %}{% assign subtotal = subtotal | plus: line %}{% endfor %}{% assign tax = subtotal | times: tax_rate | divided_by: 100.0 | round: d %}{% assign total = subtotal | plus: tax %}
\begin{tabularx}{\linewidth}{@{}m{13mm}@{\hspace{4mm}}L r@{}}
  \tikz\node[circle, fill=accent, text=white, minimum size=12mm, font=\Large\bfseries, inner sep=0pt]{ {{- buyer.name | strip | slice: 0 | upcase -}} }; &
  \parbox{\linewidth}{\raggedright{\large\bfseries\mbox{}{{ buyer.name }}}\par{\small\mbox{}{{ buyer.address }}\par\mbox{}{{ buyer.email }}}} &
  \textcolor{accent}{\fontsize{22}{24}\selectfont\bfseries PURCHASE ORDER}
\end{tabularx}

\vspace{16pt}
{\renewcommand{\arraystretch}{1.6}\arrayrulecolor{black!20}
\begin{tabularx}{\linewidth}{|L|L|L|L|}
  \hline
  \fieldname{PO no.} & \fieldname{Date} & \fieldname{Deliver by} & \fieldname{Terms} \\
  {\bfseries {{ po_number }}} & {{ date }} & {{ delivery_date }} & {{ payment_terms }} \\
  \hline
\end{tabularx}}

\vspace{14pt}
% Each block has an accent rule the full height of its text.
\begin{tabularx}{\linewidth}{@{}!{\color{accent}\vrule width 2pt}@{\hspace{8pt}}L@{\hspace{8mm}}!{\color{accent}\vrule width 2pt}@{\hspace{8pt}}L@{}}
  \fieldname{Vendor}\par\vspace{2pt}{\bfseries\mbox{}{{ vendor.name }}}\par{\small\mbox{}{{ vendor.address }}} & \fieldname{Ship to}\par\vspace{2pt}{\bfseries\mbox{}{{ ship_to.name }}}\par{\small\mbox{}{{ ship_to.address }}}
\end{tabularx}

\vspace{18pt}
% A longtable, so a long list breaks across pages with the header repeated.
% Every column has a fixed width: one compile pass gives aligned columns.
{\renewcommand{\arraystretch}{1.6}\setlength{\LTpre}{0pt}\setlength{\LTpost}{0pt}\setlength{\LTleft}{0pt}\setlength{\LTright}{0pt}\arrayrulecolor{black!20}
\begin{longtable}{|>{\raggedright\arraybackslash}p{24mm}|>{\raggedright\arraybackslash}p{\dimexpr\linewidth-6\arrayrulewidth-10\tabcolsep-24mm-12mm-22mm-22mm\relax}|>{\raggedleft\arraybackslash}p{12mm}|>{\raggedleft\arraybackslash}p{22mm}|>{\raggedleft\arraybackslash}p{22mm}|}
  \hline
  \rowcolor{accent}\textcolor{white}{\bfseries Item no.} & \textcolor{white}{\bfseries Description} & \textcolor{white}{\bfseries Qty} & \textcolor{white}{\bfseries Unit price} & \textcolor{white}{\bfseries Amount} \\
  \hline
\endhead
{% for item in items %}{% assign price = item.unit_price | round: 4 %}{% assign line = item.quantity | times: price | round: d %}  {{ item.sku }} & {{ item.description }} & {{ item.quantity | round: 2 }} & {{ price | currency: currency, 4 }} & {{ line | currency: currency }} \\
  \hline
{% endfor %}\end{longtable}}

\vspace{8pt}
\hfill{\renewcommand{\arraystretch}{1.6}\arrayrulecolor{black!20}\begin{tabular}{|l|r|}
  \hline
  Subtotal & {{ subtotal | currency: currency }} \\
  \hline
  Tax ({{ tax_rate }}\%) & {{ tax | currency: currency }} \\
  \hline
  \rowcolor{accent}\textcolor{white}{\bfseries Total} & \textcolor{white}{\bfseries {{ total | currency: currency }}} \\
  \hline
\end{tabular}}

\vfill
\begin{tabularx}{\linewidth}{@{}L@{\hspace{10mm}}L@{}}
  \fieldname{Terms and notes} & \fieldname{Authorized by} \\[2pt]
  {\small {{ notes }}} & {\small {{ authorized_by }}\par\vspace{18pt}\textcolor{muted}{\rule{0.8\linewidth}{0.5pt}}}
\end{tabularx}
\end{document}
